% BHU PYQ -- Manually transcribed & error-corrected
\documentclass[11pt,a4paper]{article}

\usepackage[a4paper,top=2.5cm,bottom=2.5cm,left=2.8cm,right=2.8cm,headheight=14pt]{geometry}
\usepackage{amsmath,amssymb}
\usepackage[T1]{fontenc}
\usepackage[utf8]{inputenc}
\usepackage{lmodern}
\usepackage{microtype}
\usepackage{fancyhdr}
\usepackage{enumitem}
\usepackage{hyperref}
\usepackage{lastpage}
\usepackage{parskip}

\hypersetup{colorlinks=true,urlcolor=black,linkcolor=black,
  pdftitle={CHB-201: Inorganic & Physical Chemistry-I (Old Course) -- BHU PYQ},pdfauthor={Banaras Hindu University}}

\pagestyle{fancy}\fancyhf{}
\renewcommand{\headrulewidth}{0.4pt}\renewcommand{\footrulewidth}{0.4pt}
\fancyhead[L]{\small\textit{Banaras Hindu University}}
\fancyhead[R]{\small\textit{CHB-201: Inorganic \& Physical Chemistry-I (Old Course)}}
\fancyfoot[C]{\small Page \thepage\ of \pageref{LastPage}}

\newlist{parts}{enumerate}{2}
\setlist[parts,1]{label=(\alph*),leftmargin=*,itemsep=5pt,topsep=3pt}
\setlist[parts,2]{label=(\roman*),leftmargin=*,itemsep=3pt,topsep=2pt}
\newcommand{\pts}[1]{\hfill\textit{[#1]}}

\begin{document}

\begin{center}
  {\large\bfseries BANARAS HINDU UNIVERSITY}\\[2pt]
  {\normalsize B.Sc. (Hons.) Semester II Examination, 2013-14 (Old Course)}\\[6pt]
  \rule{0.8\linewidth}{0.5pt}\\[6pt]
  {\large\bfseries Paper No. CHB-201: Inorganic \& Physical Chemistry-I}\\[4pt]
  {\normalsize Time: 3 Hours\hfill Maximum Marks: 70}
\end{center}
\rule{\linewidth}{0.4pt}
\small
\noindent\textit{Instructions: Answer five questions in total, including Question~1 from each section which is compulsory. Use separate answer books for Section~A and Section~B.}
\normalsize
\bigskip

\begin{center}
  {\large\bfseries SECTION A}\\[2pt]
  {\normalsize (Inorganic Chemistry-I -- Marks: 35)}
\end{center}
\rule{0.4\linewidth}{0.2pt}
\smallskip

\noindent\textbf{Question 1.} Answer all parts of the following: \pts{5 $\times$ 2 = 10}
\begin{parts}
  \item What is electrode potential? Explain the term with the help of a Born-Haber type cycle.
  \item What happens when alkali metal tetrahydroborates are treated with water?
  \item Draw the structures of Cryptand-2,2,2 and Benzo-18-crown-6.
  \item Draw the structure of basic beryllium acetate.
  \item Draw and explain the structure of $\text{XeF}_5^-$.
\end{parts}

\medskip\rule{\linewidth}{0.2pt}

\noindent\textbf{Question 2.}
\begin{parts}
  \item Derive the Born-Land\'e equation for the lattice energy of an ionic crystal. Explain all the terms involved. \pts{6}
  \item Discuss the solvation energy of ions. What are the factors that influence the solvation of ions? \pts{6}
\end{parts}

\medskip\rule{\linewidth}{0.2pt}

\noindent\textbf{Question 3.}
\begin{parts}
  \item What are silicates? How are they classified? Give the structures of orthosilicates and chain silicates. \pts{6}
  \item Write short notes on: \pts{6}
  \begin{parts}
    \item Intercalation compounds of graphite
    \item Nitrogen activation reactions
  \end{parts}
\end{parts}

\medskip\rule{\linewidth}{0.2pt}

\noindent\textbf{Question 4.}
\begin{parts}
  \item Discuss the bonding, structure, and applications of polymeric sulfur nitride ($\text{(SN)}_x$). \dots \pts{6}
  \item Draw structures, predict hybridizations, and explain the shapes of the following: (i) $\text{XeF}_6$, (ii) $\text{XeOF}_4$, (iii) $\text{XeO}_4$. \pts{6}
\end{parts}

\newpage
\begin{center}
  {\large\bfseries SECTION B}\\[2pt]
  {\normalsize (Physical Chemistry-I -- Marks: 35)}
\end{center}
\rule{0.4\linewidth}{0.2pt}
\smallskip

\noindent\textbf{Question 5.} Answer all parts of the following: \pts{5 $\times$ 2 = 10}
\begin{parts}
  \item State the law of corresponding states and write its reduced equation.
  \item Explain why the surface tension of a liquid decreases with an increase in temperature.
  \item Explain the terms: average velocity, most probable velocity, and root-mean-square velocity.
  \item State the first law of thermodynamics. Define heat capacity at constant volume ($C_v$) and constant pressure ($C_p$).
  \item Under what conditions does a real gas behave like an ideal gas?
\end{parts}

\medskip\rule{\linewidth}{0.2pt}

\noindent\textbf{Question 6.}
\begin{parts}
  \item Define the surface tension of a liquid. Discuss the stalagmometer method for the determination of surface tension. What is the influence of temperature on surface tension? \pts{8}
  \item Water rises to a height of $7.4\,\text{cm}$ in a glass capillary at $20^\circ\text{C}$. Calculate the radius of the capillary. (Surface tension of water = $72.8\,\text{dynes/cm}$, density = $1\,\text{g/cm}^3$, $g = 981\,\text{cm/s}^2$). \pts{4}
\end{parts}

\medskip\rule{\linewidth}{0.2pt}

\noindent\textbf{Question 7.}
\begin{parts}
  \item Using the Maxwell-Boltzmann distribution of molecular speeds, derive expressions for average, most probable, and root-mean-square velocities of gas molecules. \pts{6}
  \item Discuss the kinetic theory of gases. Show how the kinetic gas equation can be used to derive Dalton's law of partial pressures. \pts{6}
\end{parts}

\medskip\rule{\linewidth}{0.2pt}

\noindent\textbf{Question 8.}
\begin{parts}
  \item What is the Joule-Thomson effect? Show that the Joule-Thomson expansion of a real gas is an isenthalpic process. \pts{6}
  \item Calculate the work done in the isothermal reversible expansion of $2\,\text{moles}$ of an ideal gas from $2\,\text{L}$ to $20\,\text{L}$ at $300\,\text{K}$. \pts{6}
\end{parts}

\vfill
\rule{\linewidth}{0.4pt}
\begin{center}\small
  \href{https://bhu-exam.vercel.app/}{Click here to visit our website}\\[4pt]
  \href{https://me-aryan.vercel.app/}{Click here to visit our contributor}\\[4pt]
  \href{https://quashan-me.vercel.app/}{Click here to visit our contributor}
\end{center}

\end{document}
